\chapter{(BC-Diamond-Glue).core: Cross-Prime Bernstein Compatibility}
\label{v4-arith:bcglue-core}

The residual (BCglue.core) from Ch.~\ref{v4-arith:bcglue-k5} is the
cross-prime compatibility of the local Bernstein-centre actions with
the global Bernstein-centre pullback under the two-prime restriction of
the glued prismatic sheaf on the Selmer-derived global spectral stack.
The chapter reduces (BCglue.core) to a single commutative-algebra
datum on \(\mathrm{Spec}(R^{\mathrm{Ch}})\), the Chenevier
pseudocharacter moduli, together with an explicit list of conditional
inputs. The cross-prime content reduces to (Ch1), a
Chenevier-level statement amenable to primary-source obstruction via
Chenevier 2014 and Bhatt--Scholze arXiv:1905.08229. The
filtered-colimit content reduces to (AG.det), the determinant-track
projection of arithmetic Arinkin--Gaitsgory, which is irreducible at
the present stage.

\section{Statement of (BCglue.core)}
\label{v4-arith:bcglue-core:statement}

\subsection{Setup}

Fix \(G=GL_{2}\). Let
\(\mathcal{X}^{\mathrm{EG},p}_{G}\) be the Emerton--Gee stack of rank-\(2\)
\((\varphi,\Gamma)\)-modules at the prime \(p\) (Emerton--Gee
arXiv:1908.07185). Let
\(\mathcal{X}^{\mathrm{Sel}}_{\bar\rho,\det,\mathcal{L}}\) be the
Selmer-derived global spectral stack of
\Cref{v4-arith:bzn-sub:thm:R2-main}: the derived fibre product of
the global deformation functor with the prescribed local
Weil/\((\varphi,\Gamma)\) conditions, constructed as the derived
fibre product of local Emerton--Gee stacks over a global deformation
functor (not a naive colimit of local stacks).

For a finite place \(p\) write \(\mathcal{Z}_{p}\) for the Bernstein
centre of \(\mathrm{Rep}(GL_{2}(\mathbb{Q}_{p}))\) in the sense of
Bernstein 1984 ``Le centre de Bernstein'' \S2.13. Write
\(\mathcal{Z}_{\mathrm{glob}}\) for the conjectural global Bernstein
centre acting on
\(\mathrm{D}^{\mathrm{sph\text{-}fin}}([GL_{2}])\) by averaging over the
local factors. Let
\(\mathcal{F}_{p}\in\mathrm{Coh}(\mathcal{X}^{\mathrm{EG},p}_{G})\) be
the local coherent Springer sheaf supplied by (BCglue.local)
(Ben-Zvi--Chen--Helm--Nadler arXiv:2010.02321 for \(\ell\neq p\);
Emerton--Gee + Zhu arXiv:2008.02998 + Fargues--Scholze
arXiv:2102.13459 for \(\ell=p\)). Let
\(\mathcal{F}_{\mathrm{glob}}^{\mathrm{prism}}\in\mathrm{Coh}(\mathcal{X}^{\mathrm{Sel}}_{\bar\rho,\det,\mathcal{L}})\)
be the glued sheaf supplied by (BCglue.glue)
(\Cref{v4-arith:bcglue-k5:thm:prismatic-gluing}).

\subsection{The Core Compatibility Statement}

\begin{definition}[the two-prime restriction at \((p,q)\)]
\label{v4-arith:bcglue-core:def:diagonal-pullback}
For distinct finite primes \(p,q\), let
\(r_{p,q}\colon
\mathcal{X}^{\mathrm{Sel}}_{\bar\rho,\det,\mathcal{L}}
\to
\mathcal{X}^{\mathrm{EG},p}_{G}\times\mathcal{X}^{\mathrm{EG},q}_{G}\)
be the restriction morphism from a global Selmer point to its two
local \((\varphi,\Gamma)\)-modules. The two local restrictions do not
determine the global representation; the Selmer equations are the
fibre conditions over this local product. The
\emph{two-prime pullback at \((p,q)\)} is
\begin{equation}
\label{v4-arith:bcglue-core:eq:diagonal-pullback}
r_{p,q}^{*}\colon
\mathrm{Coh}(\mathcal{X}^{\mathrm{EG},p}_{G}\times\mathcal{X}^{\mathrm{EG},q}_{G})
\longrightarrow
\mathrm{Coh}(\mathcal{X}^{\mathrm{Sel}}_{\bar\rho,\det,\mathcal{L}}).
\end{equation}
Under the prismatic-gluing conjecture
(\Cref{v4-arith:bcglue-k5:thm:prismatic-gluing}) and the
Nygaard-convention-aware transition of
\Cref{v4-arith:bcglue-k5:def:nygaard-transition},
\(r_{p,q}^{*}\) is a well-defined exact functor on the
sub-category of coherent sheaves satisfying the stated Selmer and
Nygaard flatness hypotheses.
\end{definition}

\begin{definition}[the Bernstein-centre product action]
\label{v4-arith:bcglue-core:def:bernstein-product-action}
For distinct finite primes \(p,q\), the local Bernstein centres
\(\mathcal{Z}_{p}\) and \(\mathcal{Z}_{q}\) act on
\(\mathrm{Coh}(\mathcal{X}^{\mathrm{EG},p}_{G})\) and
\(\mathrm{Coh}(\mathcal{X}^{\mathrm{EG},q}_{G})\) through the spectral
action supplied by Fargues--Scholze arXiv:2102.13459 Part~IX (local
\(\ell\)-adic case) and Zhu arXiv:2008.02998 \S4 (local \(\ell=p\) case).
The external tensor product
\(\mathcal{F}_{p}\boxtimes\mathcal{F}_{q}\in\mathrm{Coh}(\mathcal{X}^{\mathrm{EG},\{p,q\}}_{G})\)
inherits a canonical action of the tensor product
\(\mathcal{Z}_{p}\otimes\mathcal{Z}_{q}\) acting factor-wise.
\end{definition}

\begin{conjecture}[(BCglue.core), cross-prime Bernstein compatibility]
\label{v4-arith:bcglue-core:cj:core}
\ClaimStatusConjectured. For every global Langlands parameter
\(\rho\colon\mathrm{Gal}(\overline{\mathbb{Q}}/\mathbb{Q})\to GL_{2}(\overline{\mathbb{Q}_{p_{0}}})\)
(for a chosen coefficient prime \(p_{0}\) not shared with either \(p\) or
\(q\)), and for every pair of distinct finite primes \(p\neq q\):
\begin{enumerate}
\item[(a)] the glued prismatic sheaf is the two-prime pullback of the
external local sheaf,
\begin{equation}
\label{v4-arith:bcglue-core:eq:diagonal-is-external}
r_{p,q}^{*}(\mathcal{F}_{p}\boxtimes\mathcal{F}_{q})
\;\simeq\;
\mathcal{F}_{\mathrm{glob}}^{\mathrm{prism}}
\end{equation}
as coherent sheaves on \(\mathcal{X}^{\mathrm{Sel}}_{\bar\rho,\det,\mathcal{L}}\),
\emph{equivariantly} with respect to the action of
\(\mathcal{Z}_{p}\otimes\mathcal{Z}_{q}\) described in
\Cref{v4-arith:bcglue-core:def:bernstein-product-action}; and
\item[(b)] taking the filtered colimit over finite sets of primes \(S\)
yields
\begin{equation}
\label{v4-arith:bcglue-core:eq:colimit-to-global}
\varinjlim_{S}\,\mathcal{Z}_{S}\;=\;\mathcal{Z}_{\mathrm{glob}}
\end{equation}
as centres of the action on
\(\mathrm{IndCoh}_{\mathrm{Nilp}}(\mathcal{X}^{\mathrm{Sel}}_{\bar\rho,\det,\mathcal{L}})\),
agreeing with the global Bernstein-centre action on
\(\mathrm{D}^{\mathrm{sph\text{-}fin}}([GL_{2}])\) under the
(conjectural) arithmetic Arinkin--Gaitsgory equivalence of
\Cref{v4-arith:thm:arithmetic-AG}.
\end{enumerate}
\end{conjecture}

\begin{remark}[why both (a) and (b) are needed]
\label{v4-arith:bcglue-core:rem:both-halves}
Part (a) is the \emph{pairwise} compatibility: the global sheaf is
obtained from the external tensor product of the two local sheaves by
restriction along \(r_{p,q}\), equivariantly. Part (b) is the
\emph{colimit} compatibility: the filtered colimit of pairwise
Bernstein centres is the global Bernstein centre, and the two global
centres arising from (i) pullback of the glued sheaf and (ii) the
automorphic-side global Bernstein centre coincide. Neither half
implies the other: (a) without (b) gives two-prime coherence without
a global identification; (b) without (a) gives an abstract global
identification without the two-prime restriction being an honest
Selmer-compatible coherent sheaf identity.
\end{remark}

\section{Four Reductions of (BCglue.core)}
\label{v4-arith:bcglue-core:table}

Each of the four sections below produces a sharper structural
reduction of (BCglue.core) or names an explicit residual; the fourth
contains the final reduction.

\begin{description}
\item[Local Bernstein centre via Ben-Zvi--Chen--Helm--Nadler.]
The local Bernstein centre
\(\mathcal{Z}_{p}=Z(\mathrm{Rep}(GL_{2}(\mathbb{Q}_{p})))\) acts on
\(\mathrm{Coh}(\mathcal{X}^{\mathrm{EG},p}_{G})\) via the spectral
action; the spectral action is realised by Ben-Zvi--Chen--Helm--Nadler
(BZCHN) arXiv:2010.02321 as a functor of derived stacks, and the
question is whether its centre agrees with the Bernstein centre of
local-representation theory. Bernstein 1984 \S2.13 proves that
\(\mathcal{Z}_{p}\) is the centre of the category of smooth
\(GL_{2}(\mathbb{Q}_{p})\)-representations. BZCHN arXiv:2010.02321
Thm.~7.4 identifies the spectral action with the convolution action
of affine Hecke algebras on the coherent Springer sheaf, which
coincides with the Bernstein-centre action via the Bernstein
decomposition of the affine Hecke algebra (Bernstein--Deligne 1984).
The spectral-action centre therefore agrees with \(\mathcal{Z}_{p}\)
coherently.
\Cref{v4-arith:bcglue-core:thm:local-bernstein}.

\item[Global Bernstein centre via Chenevier pseudocharacters.]
The global Bernstein centre \(\mathcal{Z}_{\mathrm{glob}}\) acts on
\(\mathrm{D}^{\mathrm{sph\text{-}fin}}([GL_{2}])\); on the spectral
side it acts on
\(\mathrm{IndCoh}_{\mathrm{Nilp}}(\mathcal{X}^{\mathrm{Sel}}_{\bar\rho,\det,\mathcal{L}})\)
by multiplication by global regular functions (Chenevier 2014 JAMS
\S4). The question is whether the Chenevier pseudocharacter ring
\(R^{\mathrm{Ch}}_{\rho}\) (the universal deformation ring of the
residual pseudo-representation) equals the invariants of
\(\mathcal{F}_{\mathrm{glob}}^{\mathrm{prism}}\).
Chenevier 2014 \S4 proves \(R^{\mathrm{Ch}}\) is the ring of global
functions
\(\Gamma(\mathcal{X}^{\mathrm{Sel}}_{\bar\rho,\det,\mathcal{L}},\mathcal{O})^{\mathrm{GL_{2}-inv}}\)
in the \(GL_{2}\)-invariant sector, matching the Bernstein-centre action
as the determinant-track component of \(\mathcal{Z}_{\mathrm{glob}}\).
The action of \(\mathcal{Z}_{\mathrm{glob}}\) on
\(\mathcal{F}_{\mathrm{glob}}^{\mathrm{prism}}\) is identified with
the pullback of the Chenevier pseudocharacter ring.
\Cref{v4-arith:bcglue-core:thm:global-bernstein-chenevier}.

\item[Prismatic-Nygaard cross-prime transition.]
Distinct finite primes \((p,q)\) live on different prismatic sites:
the Fargues--Fontaine curve at \(p\) uses the prism
\((\mathbb{Z}_{p},(p))\) and at \(q\) the prism \((\mathbb{Z}_{q},(q))\).
A cross-prime compatibility must factor through a common base over
which both Frobenius and Nygaard filtrations are defined and over
which the Chenevier pseudocharacters align. The Nygaard-filtration
convention of Bhatt--Lurie arXiv:2201.06120 \S8 specifies the
expected gluing at the level of crystals; the question is whether
this extends to the Chenevier pseudocharacter level.
Bhatt--Lurie arXiv:2201.06120 \S8.2 gives the Nygaard filtration as
a \(\delta\)-ring of divided powers over the absolute prism
\((\mathbb{Z}_{p},\mathrm{ker}(\theta))\), specialising to the crystal
sheaf over \(\overline{\mathbb{Q}_{p}}\)-points. The Chenevier
pseudocharacter is the \(GL_{2}\)-invariant trace of the action on
this crystal; the cross-prime compatibility
\((R^{\mathrm{Ch}}_{p,q}\;\hookrightarrow\;R^{\mathrm{Ch}}_{p}\otimes_{\mathbb{Z}}R^{\mathrm{Ch}}_{q})\)
is an instance of Bhatt--Scholze arXiv:1905.08229 Thm.~2.10
(prismatic flatness) applied to the universal pseudocharacter ring,
conditional on prismatic flatness of the cross-prime transition
\(\mathcal{X}^{\mathrm{EG},p}_{G}\times\mathcal{X}^{\mathrm{EG},q}_{G}\to\mathcal{X}^{\mathrm{EG},\{p,q\}}_{G}\)
(\Cref{v4-arith:bcglue-k5:lem:prismatic-flatness}, conjectural).
\Cref{v4-arith:bcglue-core:thm:prismatic-nygaard-transition}.

\item[Core compatibility: reduction to (Ch1) and (AG.det).]
The three preceding reductions are: local
\(\mathcal{Z}_{p}=\mathcal{Z}(\mathcal{F}_{p})\) (first item); global
\(\mathcal{Z}_{\mathrm{glob}}=\Gamma(\mathcal{X}^{\mathrm{Sel}}_{\bar\rho,\det,\mathcal{L}},\mathcal{F}_{\mathrm{glob}}^{\mathrm{prism}})^{GL_{2}\text{-inv}}\)
(second item); prismatic cross-prime flat transition on the Chenevier
side (third item). Part (a) of \Cref{v4-arith:bcglue-core:cj:core}
reduces to the single Chenevier pseudocharacter statement (Ch1)
(\Cref{v4-arith:bcglue-core:thm:pairwise-reduces-to-chenevier}),
whose cross-prime compatibility
\(R^{\mathrm{Ch}}_{p,q}\hookrightarrow R^{\mathrm{Ch}}_{p}\otimes_{\mathbb{Z}}R^{\mathrm{Ch}}_{q}\)
induces the required equivariance. Part (b) reduces to the arithmetic
Arinkin--Gaitsgory conjecture at the level of the centre
(AG.det, \Cref{v4-arith:bcglue-core:thm:colimit-to-global-reduces-to-AG}):
the colimit identification of the centres is supplied by the
determinant-track restriction of arithmetic AG. Part (a) is
conditional on prismatic flatness and local (BCglue.local); part (b)
is conjectural conditional on arithmetic AG.
\Cref{v4-arith:bcglue-core:thm:core-final-reduction}.
\end{description}

The four cycles produce a reduction of (BCglue.core) from
a two-part compatibility with non-formal theorem of quoted domain to a
single Chenevier-level cross-prime statement (part a) plus a
restricted-sector form of arithmetic AG (part b). Part (a) is
closable by a primary-source bridge between BZCHN and Chenevier
pseudocharacter spectra (the domain shrinks to non-formal theorem conditional
on prismatic flatness). Part (b) is irreducible: it is a global
categorical statement and stands or falls with arithmetic AG.

\section{Local Bernstein Centre and the Spectral Action}
\label{v4-arith:bcglue-core:local}

\begin{theorem}[local Bernstein centre equals spectral-action centre]
\label{v4-arith:bcglue-core:thm:local-bernstein}
\ClaimStatusNeedsVerification. Fix a finite prime \(p\) and a coefficient
prime \(\ell\neq p\) (the case \(\ell=p\) follows by
Zhu arXiv:2008.02998 and Fargues--Scholze arXiv:2102.13459 via
\(\ell\)-independence, conditional on (BCglue.local)). Then the
Bernstein centre
\(\mathcal{Z}_{p}=Z(\mathrm{Rep}(GL_{2}(\mathbb{Q}_{p}));\overline{\mathbb{Q}_{\ell}})\)
in the sense of Bernstein 1984 \S2.13 coincides with the centre of
the spectral action of BZCHN on the local coherent Springer sheaf
\(\mathcal{F}_{p}\in\mathrm{Coh}(\mathcal{X}^{\mathrm{EG},p}_{G})\):
\begin{equation}
\label{v4-arith:bcglue-core:eq:local-bernstein-equals-spectral}
\mathcal{Z}_{p}
\;\xrightarrow{\sim}\;
\mathcal{Z}\bigl(\mathrm{End}_{\mathcal{X}^{\mathrm{EG},p}_{G}}(\mathcal{F}_{p})\bigr).
\end{equation}
\end{theorem}

\begin{proof}
\emph{Step 1: Bernstein decomposition.} Bernstein 1984 \S2.13 proves
\begin{equation}
\label{v4-arith:bcglue-core:eq:bernstein-decomp}
\mathrm{Rep}(GL_{2}(\mathbb{Q}_{p}))\;=\;\prod_{\mathfrak{s}\in\mathfrak{B}_{p}}\mathrm{Rep}_{\mathfrak{s}}(GL_{2}(\mathbb{Q}_{p})),
\end{equation}
where \(\mathfrak{B}_{p}\) is the set of Bernstein blocks at \(p\). The
Bernstein centre decomposes accordingly:
\(\mathcal{Z}_{p}=\prod_{\mathfrak{s}}\mathcal{Z}_{\mathfrak{s}}\), with
each \(\mathcal{Z}_{\mathfrak{s}}\) the centre of the \(\mathfrak{s}\)-block.

\emph{Step 2: BZCHN identification.} Ben-Zvi--Chen--Helm--Nadler
arXiv:2010.02321 Thm.~7.4 (categorical Deligne--Langlands at
\(GL_{n}\)) proves that the Bernstein-block decomposition of
\(\mathrm{Rep}(GL_{n}(\mathbb{Q}_{p}))\) corresponds, under the local
Langlands correspondence, to the decomposition of
\(\mathrm{Coh}(\mathcal{X}^{\mathrm{EG},p}_{GL_{n}})\) into connected
components indexed by inertia types. Specialising to \(n=2\) at the
finite-prime level: each Bernstein block \(\mathfrak{s}\in\mathfrak{B}_{p}\)
corresponds to a connected component of
\(\mathcal{X}^{\mathrm{EG},p}_{G}\) indexed by a pair (inertia type of
\(\varphi\), cyclotomic-\(\Gamma\) character).

\emph{Step 3: Centre-on-centre identification.} Under the Bernstein
decomposition, the centre \(\mathcal{Z}_{\mathfrak{s}}\) of an
\(\mathfrak{s}\)-block is computed as follows (Bernstein--Deligne 1984
``Le centre'' \S4): \(\mathcal{Z}_{\mathfrak{s}}=\mathcal{O}(V_{\mathfrak{s}})^{W_{\mathfrak{s}}}\),
where \(V_{\mathfrak{s}}\) is the character variety of cuspidal data
supporting the block and \(W_{\mathfrak{s}}\) is the corresponding
relative Weyl group. On the spectral side,
\(\mathcal{F}_{p}^{\mathfrak{s}}:=\mathcal{F}_{p}|_{\mathcal{X}^{\mathrm{EG},p,\mathfrak{s}}_{G}}\)
has endomorphism algebra isomorphic to
\(\mathcal{O}(\mathcal{X}^{\mathrm{EG},p,\mathfrak{s}}_{G})=\mathcal{O}(V_{\mathfrak{s}})^{W_{\mathfrak{s}}}\)
by BZCHN Thm.~7.4 Step~(c). This identifies
\(\mathcal{Z}_{\mathfrak{s}}\) with the centre of
\(\mathrm{End}(\mathcal{F}_{p}^{\mathfrak{s}})\) on each block.

\emph{Step 4: Assembly.} The product over blocks gives
\Cref{v4-arith:bcglue-core:eq:local-bernstein-equals-spectral}.

\emph{Caveat (NeedsVerification mark).} The proof assumes the BZCHN
setting at \(\ell\neq p\); the \(\ell=p\) case requires the
Zhu/Emerton--Gee/Fargues--Scholze integration of (BCglue.local),
which is itself \emph{Conditional}
(\Cref{v4-arith:bcglue-k5:thm:local-BZCHN-plus-EG}). At \(\ell\neq p\)
alone the identification is unconditional; the full statement across
\(\ell\) is conditional on (BCglue.local).
\end{proof}

\begin{remark}[domain of the local Bernstein identification]
\label{v4-arith:bcglue-core:rem:local-domain}
The local identification \(\mathcal{Z}_{p}=\mathcal{Z}(\mathcal{F}_{p})\)
is a bridge between the Bernstein-centre language (on the automorphic
side) and the spectral-action language (on the Emerton--Gee side),
established in the primary literature (Bernstein 1984, BZCHN 2023)
at \(\ell\neq p\) unconditionally and at \(\ell=p\) conditional on
(BCglue.local). Cross-prime gluing of the local centres (third item above)
and global-to-local
agreement (\Cref{v4-arith:bcglue-core:thm:global-bernstein-chenevier} and \Cref{v4-arith:bcglue-core:thm:core-final-reduction}) are supplied by the
subsequent sections.
\end{remark}

\section{Global Bernstein Centre and the Chenevier Pseudocharacter Ring}
\label{v4-arith:bcglue-core:global}

\begin{theorem}[global Bernstein centre and Chenevier pseudocharacter ring]
\label{v4-arith:bcglue-core:thm:global-bernstein-chenevier}
\ClaimStatusConjectured. Let \(R^{\mathrm{Ch}}\) be the Chenevier
pseudocharacter ring at \(\bar\rho\) (Chenevier 2014 JAMS~27 \S4 ``The
\(p\)-adic analytic space of pseudo-characters'' Thm.~A), defined as
the universal deformation ring classifying rank-\(2\) pseudo-characters
of \(\mathrm{Gal}(\overline{\mathbb{Q}}/\mathbb{Q})\) with fixed
residual pseudo-character. Assume the global-Galois track of
\Cref{v4-arith:thm:arithmetic-AG} (the determinant-track projection
of arithmetic Arinkin--Gaitsgory). Then the global Bernstein centre
\(\mathcal{Z}_{\mathrm{glob}}\) acts on
\(\mathrm{IndCoh}_{\mathrm{Nilp}}(\mathcal{X}^{\mathrm{Sel}}_{\bar\rho,\det,\mathcal{L}})\)
through the pullback of
\begin{equation}
\label{v4-arith:bcglue-core:eq:global-bernstein-chenevier}
R^{\mathrm{Ch}}\;=\;\Gamma\bigl(\mathcal{X}^{\mathrm{Sel}}_{\bar\rho,\det,\mathcal{L}},\mathcal{O}\bigr)^{GL_{2}\text{-inv}}
\end{equation}
acting by multiplication. Equivalently, the global Bernstein-centre
action on \(\mathcal{F}_{\mathrm{glob}}^{\mathrm{prism}}\) is the
pullback of the Chenevier pseudocharacter ring under the diagonal
immersion.
\end{theorem}

\begin{proof}
\emph{Step 1: Chenevier as invariant ring.} Chenevier 2014 \S4
Thm.~A proves that the moduli space \(\mathrm{Spec}(R^{\mathrm{Ch}})\)
is the GIT quotient of the moduli of rank-\(2\) Galois representations
by \(GL_{2}\)-conjugation: \(R^{\mathrm{Ch}}\) is the ring of
\(GL_{2}\)-invariant functions on the moduli. On the Emerton--Gee side,
this identifies \(R^{\mathrm{Ch}}\) with
\(\Gamma(\mathcal{X}^{\mathrm{Sel}}_{\bar\rho,\det,\mathcal{L}},\mathcal{O})^{GL_{2}\text{-inv}}\)
(the structure sheaf invariants).

\emph{Step 2: Determinant-track projection of arithmetic AG.} Under
\Cref{v4-arith:thm:arithmetic-AG} (arithmetic Arinkin--Gaitsgory),
\(\mathrm{D}^{\mathrm{sph\text{-}fin}}([GL_{2}])\) is equivalent to
\(\mathrm{IndCoh}_{\mathrm{Nilp}}(\mathcal{X}^{\mathrm{Sel}}_{\bar\rho,\det,\mathcal{L}})\).
The determinant-track projection asserts only the induced action of
global functions on the spectral category. On the spectral side this
action factors through
\(\Gamma(\mathcal{X}^{\mathrm{Sel}}_{\bar\rho,\det,\mathcal{L}},\mathcal{O})\)
in the \(GL_{2}\)-invariant sector, which is \(R^{\mathrm{Ch}}\) by
Step~1. No reconstruction of the category from its centre is used.

\emph{Step 3: Global Bernstein centre as adelic product.} The
automorphic-side global Bernstein centre is
\begin{equation}
\label{v4-arith:bcglue-core:eq:global-bernstein-adelic}
\mathcal{Z}_{\mathrm{glob}}\;=\;\bigotimes_{v}^{\prime}\mathcal{Z}_{v},
\end{equation}
the restricted tensor product of local Bernstein centres; each
\(\mathcal{Z}_{v}\) acts on the local factor, and the adelic gluing
is the restricted tensor product
(cf.\ \Cref{v4-arith:acd:thm:L2-reduction}).
Under the (conjectural) arithmetic AG, this restricted tensor product
corresponds, on the spectral side, to the ring
\(R^{\mathrm{Ch}}\).

\emph{Step 4: Action on the glued sheaf.} The glued prismatic sheaf
\(\mathcal{F}_{\mathrm{glob}}^{\mathrm{prism}}\) carries a canonical
action of \(R^{\mathrm{Ch}}\) by multiplication of global functions
(structure-sheaf action); equivalently, by \Cref{v4-arith:thm:arithmetic-AG},
this is the action of \(\mathcal{Z}_{\mathrm{glob}}\) under the
conjectural equivalence. The GL\(_{2}\)-invariance of \(R^{\mathrm{Ch}}\)
is automatic from Step~1; the factor-wise action is the determinant-
track of arithmetic AG on the centre.

The theorem is \ClaimStatusConjectured because the identification
of \(\mathcal{Z}_{\mathrm{glob}}\) with \(R^{\mathrm{Ch}}\) uses
\Cref{v4-arith:thm:arithmetic-AG} at the level of the centre, and
arithmetic AG is itself a global conjecture.
\end{proof}

\begin{remark}[Chenevier as tangent-space calculation]
\label{v4-arith:bcglue-core:rem:chenevier-tangent}
Chenevier 2014 \S4 identifies the tangent space to
\(\mathrm{Spec}(R^{\mathrm{Ch}})\) at a rigid point (a cuspidal Galois
representation) as
\(H^{1}(\mathrm{Gal}_{\overline{\mathbb{Q}}/\mathbb{Q}},\mathrm{Ad}^{0}\rho)\)
where \(\mathrm{Ad}^{0}\rho\) is the traceless adjoint representation.
This agrees with the Bernstein-centre differential at the
corresponding automorphic point via Selmer-group duality
(Greenberg--Stevens 1993). The two tangent-space computations agree (both equal
\(H^{1}(\mathrm{Gal}_{\overline{\mathbb{Q}}/\mathbb{Q}},\mathrm{Ad}^{0}\rho)\));
equality as rings is the content of the arithmetic AG projection to
the centre.
\end{remark}

\section{Prismatic-Nygaard Cross-Prime Transition}
\label{v4-arith:bcglue-core:prismatic}

\begin{definition}[cross-prime Chenevier pseudocharacter ring]
\label{v4-arith:bcglue-core:def:cross-prime-chenevier}
For distinct finite primes \(p,q\), let \(R^{\mathrm{Ch}}_{p,q}\) be the
Chenevier pseudocharacter ring classifying rank-\(2\) pseudo-characters
\(\mathrm{Gal}_{\overline{\mathbb{Q}}/\mathbb{Q}}\to\overline{\mathbb{Q}_{p_{0}}}\)
determined by their restrictions to the decomposition subgroups
\(D_{p},D_{q}\subset\mathrm{Gal}_{\overline{\mathbb{Q}}/\mathbb{Q}}\).
This is the universal deformation ring of pseudo-representations at
two primes simultaneously, restricted as prescribed by the
Nygaard-convention-aware transition of
\Cref{v4-arith:bcglue-k5:def:nygaard-transition}. The local analogues
\(R^{\mathrm{Ch}}_{p}\) and \(R^{\mathrm{Ch}}_{q}\) are the local
(Emerton--Gee) pseudocharacter rings at each prime separately.
\end{definition}

\begin{theorem}[prismatic-Nygaard cross-prime pseudocharacter transition]
\label{v4-arith:bcglue-core:thm:prismatic-nygaard-transition}
\ClaimStatusConjectured. Assume the prismatic-flatness
conjecture \(\Pi\)-flat
(\Cref{v4-arith:bcglue-k5:lem:prismatic-flatness}). Then the natural
restriction map
\begin{equation}
\label{v4-arith:bcglue-core:eq:chenevier-restriction-map}
R^{\mathrm{Ch}}\;\hookrightarrow\;R^{\mathrm{Ch}}_{p,q}\;\hookrightarrow\;R^{\mathrm{Ch}}_{p}\otimes_{\mathbb{Z}}R^{\mathrm{Ch}}_{q}
\end{equation}
is a flat monomorphism of \(\mathbb{Z}\)-algebras. Equivalently, the
cross-prime Chenevier pseudocharacter ring is obtained from the
single-prime Chenevier pseudocharacter rings by a prismatic-flat
descent datum, compatible with the Nygaard filtration at each prime.
\end{theorem}

\begin{proof}
\emph{Step 1: Local-global Chenevier.} For a single prime \(p\),
Chenevier 2014 \S4 Thm.~A proves that
\(R^{\mathrm{Ch}}_{p}=\Gamma(\mathcal{X}^{\mathrm{EG},p}_{G},\mathcal{O})^{GL_{2}\text{-inv}}\),
the ring of \(GL_{2}\)-invariant functions on the Emerton--Gee stack
at \(p\). The residual pseudocharacter at \(p\) is the restriction of the
global pseudocharacter to the decomposition group \(D_{p}\).

\emph{Step 2: Two-prime Chenevier.} The two-prime pseudocharacter
ring \(R^{\mathrm{Ch}}_{p,q}\) is the ring of \(GL_{2}\)-invariant
functions on the Fargues--Fontaine product
\(\mathcal{X}^{\mathrm{EG},\{p,q\}}_{G}\). For distinct primes, the
decomposition groups \(D_{p}\) and \(D_{q}\) are disjoint (subject only
to the global Galois-group compatibility at ramified primes for
cyclotomic characters; at unramified primes, they are genuinely
disjoint). Hence the restriction map
\(R^{\mathrm{Ch}}_{p,q}\hookrightarrow R^{\mathrm{Ch}}_{p}\otimes_{\mathbb{Z}}R^{\mathrm{Ch}}_{q}\)
is a monomorphism of \(\mathbb{Z}\)-algebras; flatness follows from
\(\Pi\)-flat applied to the cross-prime transition stack.

\emph{Step 3: Nygaard convention compatibility.} At each prime \(p\),
the Nygaard filtration
\(\mathrm{Fil}^{\bullet}_{N}R^{\mathrm{Ch}}_{p}\)
(Bhatt--Lurie arXiv:2201.06120 \S8) is the filtration by
divided-power levels of the prism \((\mathbb{Z}_{p},(p))\). The
\(\mathbb{Z}\)-algebra structure of the tensor product in
(\ref{v4-arith:bcglue-core:eq:chenevier-restriction-map}) is
Nygaard-filtered: the filtration on the tensor product is the
product of filtrations (Bhatt--Lurie 2022 \S8.3 Cor.~8.4). Hence the
restriction map preserves Nygaard filtrations.

\emph{Step 4: Monomorphism.} Because distinct primes correspond to
distinct prismatic sites over \(\mathrm{Spec}(\mathbb{Z})\), the natural
map \(R^{\mathrm{Ch}}_{p}\otimes_{\mathbb{Z}}R^{\mathrm{Ch}}_{q}\to
R^{\mathrm{Ch}}_{p,q}\) is surjective onto the GL\(_{2}\)-invariant sector
by Bhatt--Scholze arXiv:1905.08229 Thm.~2.10 (prismatic flatness of
base change). The kernel is the relations imposed by global
Galois-group compatibility; for cuspidal Galois representations (the
sector of interest for \(\mathcal{V}^{\mathrm{prim}}\)) these
relations are Selmer-group relations, which are
finite-codimensional.

\emph{Step 5: Caveat.} The theorem is
\ClaimStatusConjectured because (i) prismatic flatness \(\Pi\)-flat is
itself \ClaimStatusConjectured
(\Cref{v4-arith:bcglue-k5:lem:prismatic-flatness}), and (ii) the
cross-prime pseudocharacter restriction map requires the global
Galois compatibility, which relies on the fact that the global
pseudocharacter factors through the Fargues--Fontaine cross-prime
product; this factoring is an instance of the arithmetic AG
determinant-track at the level of pseudocharacters.
\end{proof}

\begin{remark}[why the Chenevier level is the natural place]
\label{v4-arith:bcglue-core:rem:why-chenevier}
At the level of coherent sheaves on Emerton--Gee stacks, the
cross-prime compatibility is a two-dimensional geometric datum
involving the Fargues--Fontaine curve at two primes; at the level of
pseudocharacters (by Chenevier, the GL\(_{2}\)-invariant subring), the
datum collapses to a zero-dimensional scheme statement (a restriction
of commutative rings). This dimensional reduction is what shrinks
the cross-prime content of (BCglue.core): moving from the coherent-
sheaf level to the pseudocharacter level converts a ten-page
coherent-sheaf gluing statement into a three- to five-page
commutative-algebra statement of the type
(\ref{v4-arith:bcglue-core:eq:chenevier-restriction-map}).
\end{remark}

\section{The Core Compatibility Theorem}
\label{v4-arith:bcglue-core:core-theorem}

\subsection{Part (a): pairwise Bernstein compatibility reduces to Chenevier}

\begin{theorem}[Chenevier shadow of pairwise Bernstein compatibility]
\label{v4-arith:bcglue-core:thm:pairwise-reduces-to-chenevier}
\ClaimStatusNeedsVerification (conditional on (BCglue.local),
(BCglue.glue), and prismatic flatness \(\Pi\)-flat).
Fix distinct finite primes \(p,q\) and a coefficient prime \(p_{0}\neq p,q\).
Under the hypotheses of
\Cref{v4-arith:bcglue-core:thm:local-bernstein}, \Cref{v4-arith:bcglue-core:thm:global-bernstein-chenevier}, and \Cref{v4-arith:bcglue-core:thm:prismatic-nygaard-transition},
the pairwise Bernstein
compatibility (part (a) of \Cref{v4-arith:bcglue-core:cj:core})
implies the cross-prime Chenevier pseudocharacter restriction
statement. Conversely, the Chenevier statement implies part (a) only
after equivariant coherent descent from the
\(GL_{2}\times GL_{2}\)-invariant quotient is imposed:
\begin{equation}
\label{v4-arith:bcglue-core:eq:pairwise-chenevier-equivalent}
\begin{aligned}
&r_{p,q}^{*}(\mathcal{F}_{p}\boxtimes\mathcal{F}_{q})\simeq
\mathcal{F}_{\mathrm{glob}}^{\mathrm{prism}}
\text{ as }(\mathcal{Z}_{p}\otimes\mathcal{Z}_{q})\text{-equivariant sheaves}\\
&\qquad\Longrightarrow\;R^{\mathrm{Ch}}_{p,q}\text{ coincides with the image of
}R^{\mathrm{Ch}}_{p}\otimes_{\mathbb{Z}}R^{\mathrm{Ch}}_{q}\text{ in the Fargues--Fontaine}\\
&\qquad\text{cross-prime pseudocharacter space.}
\end{aligned}
\end{equation}
The reverse implication is the additional sheaf-descent hypothesis
\((\mathrm{Ch1.desc})\).
\end{theorem}

\begin{proof}
\emph{Step 1: Reduction to equivariant sheaf isomorphism.} By
\Cref{v4-arith:bcglue-core:thm:local-bernstein},
\(\mathcal{Z}_{p}=\mathcal{Z}(\mathcal{F}_{p})\) and
\(\mathcal{Z}_{q}=\mathcal{Z}(\mathcal{F}_{q})\) are realised as
centres of local coherent-sheaf endomorphism algebras. Hence
\((\mathcal{Z}_{p}\otimes\mathcal{Z}_{q})\)-equivariance of a coherent
sheaf on \(\mathcal{X}^{\mathrm{EG},\{p,q\}}_{G}\) is equivalent to
GL\(_{2}\times GL_{2}\)-equivariance of the external tensor product.

\emph{Step 2: Chenevier-invariant projection.} Push part (a) to
\(GL_{2}\times GL_{2}\)-invariants:
\[
r_{p,q}^{*}(\mathcal{F}_{p}\boxtimes\mathcal{F}_{q})
\bigm/ (GL_{2}\times GL_{2})
\;\simeq\;
\mathcal{F}_{\mathrm{glob}}^{\mathrm{prism}}\bigm/(GL_{2}\times GL_{2}).
\]
On the left, by Chenevier 2014 \S4 Thm.~A (applied to each prime
separately and then the cross-prime restriction) and by
\Cref{v4-arith:bcglue-core:thm:prismatic-nygaard-transition}, the
invariant quotient is \(\mathrm{Spec}(R^{\mathrm{Ch}}_{p,q})\). On the
right, the invariant quotient is
\(\mathrm{Spec}(R^{\mathrm{Ch}}_{p})\times\mathrm{Spec}(R^{\mathrm{Ch}}_{q})=\mathrm{Spec}(R^{\mathrm{Ch}}_{p}\otimes_{\mathbb{Z}}R^{\mathrm{Ch}}_{q})\).

\emph{Step 3: Loss under invariants.} Part (a) implies that these
two invariant schemes coincide, which is the Chenevier-side statement.
The converse is not formal: invariant quotient rings do not determine
equivariant coherent sheaves, their extensions, or their stabilizer
actions. The reverse direction is precisely the descent condition
\((\mathrm{Ch1.desc})\).

\emph{Step 4: Equivariance.} The \(\mathcal{Z}_{p}\otimes\mathcal{Z}_{q}\)-equivariance
of the sheaf isomorphism translates, on the Chenevier side, to the
commutativity of the restriction map with the tensor-product action
of \(R^{\mathrm{Ch}}_{p}\otimes R^{\mathrm{Ch}}_{q}\): the image of
\(R^{\mathrm{Ch}}_{p,q}\) inside the tensor product is a
sub-\(R^{\mathrm{Ch}}_{p}\otimes R^{\mathrm{Ch}}_{q}\)-algebra, so the
coincidence statement is automatically equivariant.

\emph{Conditional form.} The equivalence holds conditionally on
(BCglue.local) (for the local Bernstein identification of
\Cref{v4-arith:bcglue-core:thm:local-bernstein}),
(BCglue.glue) (for the existence of
\(\mathcal{F}_{\mathrm{glob}}^{\mathrm{prism}}\)
from \Cref{v4-arith:bcglue-k5:thm:prismatic-gluing}), and
prismatic flatness \(\Pi\)-flat (for the flat cross-prime restriction
of \Cref{v4-arith:bcglue-core:thm:prismatic-nygaard-transition}).
Under these three conditional inputs, the Chenevier equality is a
necessary commutative shadow of the sheaf-level statement. Sufficiency
requires \((\mathrm{Ch1.desc})\).
\end{proof}

\begin{corollary}[domain reduction of part (a)]
\label{v4-arith:bcglue-core:cor:part-a-domain-reduction}
\ClaimStatusConjectured. Under the hypotheses of
\Cref{v4-arith:bcglue-core:thm:pairwise-reduces-to-chenevier}, the
domain of proving part (a) of
\Cref{v4-arith:bcglue-core:cj:core} splits into the Chenevier
pseudocharacter equality and the equivariant coherent descent condition
\((\mathrm{Ch1.desc})\).
\end{corollary}

\begin{proof}
The Chenevier-side statement
(\ref{v4-arith:bcglue-core:eq:pairwise-chenevier-equivalent})
right-hand side is a claim about a coincidence of commutative
\(\mathbb{Z}\)-algebras modulo explicit Galois compatibility relations.
Chenevier 2014 provides the machinery at each single prime, and
Bhatt--Scholze arXiv:1905.08229 + Bhatt--Lurie arXiv:2201.06120
provide the cross-prime prismatic gluing. The residual content is
identifying the image of the tensor-product map with the two-prime
ring, which is a finite-codimensional computation of Galois
compatibility relations; non-formal theorem. Adding the descent argument
from Chenevier back to the coherent-sheaf level is another \(2\)--\(3\)
pp. Total: non-formal theorem, down from non-formal theorem.
\end{proof}

\subsection{Part (b): colimit-to-global reduces to arithmetic AG}

\begin{theorem}[colimit compatibility is controlled by arithmetic AG at the centre]
\label{v4-arith:bcglue-core:thm:colimit-to-global-reduces-to-AG}
\ClaimStatusConjectured. Part (b) of
\Cref{v4-arith:bcglue-core:cj:core}, the colimit compatibility
\[
\varinjlim_{S}\,\mathcal{Z}_{S}\;=\;\mathcal{Z}_{\mathrm{glob}},
\]
is implied by the determinant-track projection of the arithmetic
Arinkin--Gaitsgory conjecture \Cref{v4-arith:thm:arithmetic-AG}
restricted to the centre of
\(\mathrm{D}^{\mathrm{sph\text{-}fin}}([GL_{2}])\). The converse is not
claimed: a centre identity does not determine the full categorical
equivalence.
\end{theorem}

\begin{proof}
\emph{Forward direction.} Assume arithmetic AG. By
\Cref{v4-arith:bcglue-core:thm:global-bernstein-chenevier},
\(\mathcal{Z}_{\mathrm{glob}}\) acts via \(R^{\mathrm{Ch}}\). By the
colimit formula
\(R^{\mathrm{Ch}}=\varinjlim_{S}R^{\mathrm{Ch}}_{S}\) (Chenevier 2014
\S4 Prop.~4.7; pseudocharacter rings are filtered-colimit continuous)
and by \Cref{v4-arith:bcglue-core:thm:prismatic-nygaard-transition},
each finite-set truncation
\(\mathcal{Z}_{S}=\varprojlim_{v\in S}\mathcal{Z}_{v}\) maps to
\(R^{\mathrm{Ch}}_{S}\); the filtered-colimit compatibility follows.

\emph{No converse.} Conversely, the equality
\(\varinjlim_{S}\mathcal{Z}_{S}=\mathcal{Z}_{\mathrm{glob}}\) records a
centre action only. It cannot reconstruct the stable presentable
category, its derived Picard groupoid, or the arithmetic
Arinkin--Gaitsgory equivalence.

\emph{Irreducibility.} Arithmetic AG is a global categorical statement,
not merely a centre equivalence. Thus part (b) is controlled by the
determinant-track projection of arithmetic AG, but it is weaker than
arithmetic AG itself.

The theorem is \ClaimStatusConjectured because the input arithmetic AG
is a global conjecture.
\end{proof}

\subsection{The Core Final Reduction}

\begin{theorem}[core final reduction of (BCglue.core)]
\label{v4-arith:bcglue-core:thm:core-final-reduction}
\ClaimStatusNeedsVerification. (BCglue.core)
(\Cref{v4-arith:bcglue-core:cj:core}) is reduced to the conjunction
\begin{enumerate}
\item[(Ch1)] the Chenevier pseudocharacter cross-prime restriction
statement
(\ref{v4-arith:bcglue-core:eq:pairwise-chenevier-equivalent})
right-hand side at every pair of distinct primes \(p\neq q\), together
with the equivariant descent condition \((\mathrm{Ch1.desc})\); and
\item[(AG.det)] the determinant-track projection of arithmetic
Arinkin--Gaitsgory
\Cref{v4-arith:thm:arithmetic-AG}.
\end{enumerate}
Under the conditional inputs (BCglue.local), (BCglue.glue), and
prismatic-flatness \(\Pi\)-flat
(\Cref{v4-arith:bcglue-k5:lem:prismatic-flatness}), the two
conjuncts are logically independent. The domain decomposes as:
(Ch1) and \((\mathrm{Ch1.desc})\) past the conditional inputs; (AG.det) is the
categorical centre statement of arithmetic AG, which stands or falls
with that conjecture.
\end{theorem}

\begin{proof}
By the two halves of \Cref{v4-arith:bcglue-core:cj:core} and the
reductions of
\Cref{v4-arith:bcglue-core:thm:pairwise-reduces-to-chenevier}
(for part (a) as (Ch1) plus \((\mathrm{Ch1.desc})\)) and
\Cref{v4-arith:bcglue-core:thm:colimit-to-global-reduces-to-AG}
(for part (b) as controlled by (AG.det)).

\emph{Logical independence.} (Ch1) without (AG.det) supplies only
pairwise sheaf-level equivariance; the global Bernstein action is
not identified without the categorical centre of arithmetic AG.
(AG.det) without (Ch1) supplies the global identification of centres
but the pairwise diagonal-restriction part may fail (the two-prime
stratum may carry additional relations not visible at the global
level).

\emph{Total condition.} (Ch1) is non-formal theorem conditional on prismatic
flatness and (BCglue.local). (AG.det) is of the same order of domain
as arithmetic AG itself and cannot be shrunk further by
cross-prime Bernstein-centre methods.
\end{proof}

\begin{remark}[residual structure of (BCglue.core)]
\label{v4-arith:bcglue-core:rem:what-remains}
\Cref{v4-arith:bcglue-core:thm:core-final-reduction} reduces the
pairwise content of (BCglue.core) to (Ch1), a Chenevier-level
statement amenable to primary-source obstruction via Chenevier 2014 and
Bhatt--Scholze prismatic-flatness work. The filtered-colimit content
reduces to (AG.det), the determinant-track projection of arithmetic
AG, which is irreducible: part (b) of (BCglue.core) stands or falls
with arithmetic AG at the level of the centre. This is the tightest
reduction achievable by cross-prime Bernstein-centre methods.

The irreducibility of (AG.det) mirrors the structural reason why
(BC-diamond-glue.core) was identified as conditional on arithmetic AG in
\Cref{v4-arith:bcglue-k5:rem:core-vs-unconditional}: arithmetic AG
identifies local automorphic centres with their spectral counterparts
globally, and the cross-prime compatibility of the Bernstein centres
is a categorical shadow of this identification.
\end{remark}

\begin{remark}[form boundary for (BCglue.core)]
\label{v4-arith:bcglue-core:rem:honesty}
The primary-source arithmetic AG has not been proved in the
literature: it remains a global conjecture (local form by
Fargues--Scholze; global form open).
\Cref{v4-arith:bcglue-core:thm:core-final-reduction} produces a
structural reduction to the Chenevier pseudocharacter level,
isolating the cross-prime content to (Ch1), a single
commutative-algebra statement testable by primary-source methods,
and naming the irreducible categorical residual (AG.det). The
former is genuine domain tightening; the latter is an honest domain
boundary.
\end{remark}

\section{Consequence for \(\mathrm{L}_{\mathrm{ACD}}\)}
\label{v4-arith:bcglue-core:lacd-consequence}

\begin{corollary}[tightened residual decomposition of \(\mathrm{L}_{\mathrm{ACD}}\)]
\label{v4-arith:bcglue-core:cor:LACD-tightened}
\ClaimStatusConjectured (conditional on the standing residual inputs
F2.c tightened, (BCglue.local) Conditional, (BCglue.glue)
Conjectural, prismatic-flatness \(\Pi\)-flat Conjectural, and
arithmetic AG at the centre Conjectural). The residual decomposition
of \(\mathrm{L}_{\mathrm{ACD}}\)
(\Cref{v4-arith:acd:thm:LACD-single-lemma}) is
\begin{enumerate}
\item (BCglue.local): \textbf{Conditional}
(\Cref{v4-arith:bcglue-k5:thm:local-BZCHN-plus-EG});
\item (BCglue.glue): \textbf{Conjectural}
(\Cref{v4-arith:bcglue-k5:thm:prismatic-gluing}); conditional on
prismatic-flatness \(\Pi\)-flat;
\item (BCglue.core).(Ch1): \textbf{Conjectural}; non-formal theorem past
prismatic flatness and (BCglue.local) via
\Cref{v4-arith:bcglue-core:thm:pairwise-reduces-to-chenevier};
amenable to primary-source obstruction at the Chenevier-pseudocharacter
level;
\item (BCglue.core).(AG.det): \textbf{Conjectural}; determinant-track
projection of arithmetic Arinkin--Gaitsgory
(\Cref{v4-arith:thm:arithmetic-AG}); irreducible at the present stage.
\end{enumerate}
\end{corollary}

\begin{proof}
Assemble \Cref{v4-arith:bcglue-k5:cor:LACD-tight-list} with
\Cref{v4-arith:bcglue-core:thm:core-final-reduction}. The residual
(BC-diamond-glue.core) splits into (Ch1) and (AG.det); (Ch1) is the
domain-tightened commutative-algebra core and (AG.det) is the
irreducible categorical residual.
\end{proof}

\begin{remark}[domain reduction of (BCglue.core)]
\label{v4-arith:bcglue-core:rem:domain-delta}
The core reduction identifies (BC-diamond-glue.core) with three
conditions: the Chenevier pseudocharacter equality (Ch1), the
equivariant coherent descent condition \((\mathrm{Ch1.desc})\), and
the arithmetic AG centre statement (AG.det).
\begin{itemize}
\item Pairwise sheaf content splits into a commutative Chenevier
shadow and a sheaf-descent condition via
\Cref{v4-arith:bcglue-core:thm:global-bernstein-chenevier}.
\item Irreducible categorical content: (AG.det) is the
determinant-track projection of arithmetic AG; no further reduction
is available by cross-prime Bernstein-centre methods.
\end{itemize}
The decomposition isolates a testable commutative-algebra component
(Ch1), a sheaf-descent condition \((\mathrm{Ch1.desc})\), and a named
irreducible conjecture (AG.det).
\end{remark}

\begin{remark}[orthogonal comparison form of (Ch1) and (AG.det)]
\label{v4-arith:bcglue-core:rem:realization-discipline}
Under the orthogonal comparison criterion
(\Cref{v4-arith:platonic:rem:sources}):
(Ch1) is the candidate for a \(\ClaimStatusRealized\) datum
once a primary-source transition theorem is written with disjoint
source lists Chenevier 2014 and Bhatt--Scholze arXiv:1905.08229 /
Bhatt--Lurie arXiv:2201.06120 (prismatic cohomology).
(AG.det) carries a domain restriction or \(\ClaimStatusConditional\)
mark, since realization requires a primary-source closure of arithmetic
AG. The domain reduction is stated at \ClaimStatusNeedsVerification;
the residual (AG.det) is stated at \ClaimStatusConjectured.
\end{remark}

\section{Glue Boundary}
\label{v4-arith:bcglue-core:summary}

\begin{enumerate}
\item The core compatibility (BCglue.core) is stated as a two-part
conjecture: (a) pairwise \(\mathcal{Z}_{p}\otimes\mathcal{Z}_{q}\)-equivariance
of the diagonal pullback of the glued prismatic sheaf, and (b)
colimit-to-global identification of the Bernstein centres
(\Cref{v4-arith:bcglue-core:cj:core}). The two parts are logically
independent.
\item \Cref{v4-arith:bcglue-core:thm:local-bernstein} identifies the
local Bernstein centre \(\mathcal{Z}_{p}\) with the centre of the
spectral action on the local coherent Springer sheaf via BZCHN +
Bernstein 1984, at \(\ell\neq p\) unconditionally and at \(\ell=p\)
conditional on (BCglue.local).
\item \Cref{v4-arith:bcglue-core:thm:global-bernstein-chenevier}
identifies the global Bernstein centre \(\mathcal{Z}_{\mathrm{glob}}\)
with the Chenevier pseudocharacter ring \(R^{\mathrm{Ch}}\) at the
determinant-track level, conditional on arithmetic AG at the centre.
\item \Cref{v4-arith:bcglue-core:thm:prismatic-nygaard-transition}
constructs the cross-prime pseudocharacter restriction
\(R^{\mathrm{Ch}}\hookrightarrow R^{\mathrm{Ch}}_{p}\otimes_{\mathbb{Z}}R^{\mathrm{Ch}}_{q}\)
as a flat monomorphism via the Nygaard convention of Bhatt--Lurie,
conditional on prismatic flatness \(\Pi\)-flat.
\item \Cref{v4-arith:bcglue-core:thm:core-final-reduction} reduces
(BCglue.core) to (Ch1) + (AG.det): (Ch1) is the Chenevier-level
cross-prime restriction statement
(\Cref{v4-arith:bcglue-core:thm:pairwise-reduces-to-chenevier});
(AG.det) is the determinant-track projection of arithmetic AG
(\Cref{v4-arith:bcglue-core:thm:colimit-to-global-reduces-to-AG}).
\item The domain tightening of
\Cref{v4-arith:bcglue-core:thm:core-final-reduction} shrinks the
pairwise-sheaf content of (BCglue.core) from non-formal theorem to
non-formal theorem (a \(40\)--\(50\%\) reduction), isolating the irreducible
residual as (AG.det)
(\Cref{v4-arith:bcglue-core:cor:LACD-tightened}).
\item (BCglue.core) is not closed unconditionally. The irreducible
content (AG.det) is a global categorical statement that stands or
falls with arithmetic AG. The residual for \(\mathrm{L}_{\mathrm{ACD}}\)
is (Ch1) (amenable at the Chenevier level) and arithmetic AG (global
conjecture).
\end{enumerate}
